\documentclass[11pt]{article}
\usepackage[margin=1in]{geometry}
\usepackage{amsmath,amssymb,amsthm,mathtools}
\usepackage{booktabs}
\usepackage{enumitem}
\usepackage{array}
\usepackage{longtable}
\usepackage[hidelinks]{hyperref}

\newtheorem{theorem}{Theorem}[section]
\newtheorem{lemma}[theorem]{Lemma}
\newtheorem{fact}[theorem]{Fact}
\newtheorem{corollary}[theorem]{Corollary}
\newtheorem{proposition}[theorem]{Proposition}
\theoremstyle{definition}
\newtheorem{definition}[theorem]{Definition}
\theoremstyle{remark}
\newtheorem{remark}[theorem]{Remark}

\newcommand{\Fh}{\widehat{F}}
\newcommand{\fh}{\widehat{f}}
\newcommand{\E}{\mathbb{E}}
\newcommand{\Prob}{\mathbb{P}}
\newcommand{\ceil}[1]{\left\lceil #1 \right\rceil}
\newcommand{\floor}[1]{\left\lfloor #1 \right\rfloor}
\newcommand{\code}[1]{\texttt{#1}}
\newcommand{\xmin}{x_{\min}}
\newcommand{\xmax}{x_{\max}}
\newcommand{\G}{\mathcal{G}}

\title{Why \code{distribution\_sort} Runs in Expected Linear Time}
\author{Analysis of \code{Linear Time Sort/distribution\_sort.py}}
\date{October 2026}

\begin{document}
\maketitle

\begin{abstract}
\code{distribution\_sort} fits a model $\Fh$ of the cumulative distribution
function (CDF) of its input, drops each element $x$ into bucket
$\floor{n\,\Fh(x)}$, and sorts each bucket. We prove that it always returns
the input in sorted order; that on every input it takes $O(n\log n)$ time;
and that when the values are independent draws from a continuous
distribution, its expected time is
$(34+5\chi^2)\,n+O(\sqrt n)$, where $\chi^2$ is the $\chi^2$-distance
between the data's density and the model's density. In particular the
expected time is linear whenever the model fits, and even when it fits
only roughly. We also prove the limit of the method: if one value repeats
a constant fraction of the time, no CDF model can avoid
$\Omega(n\log n)$ time, which is what the benchmark on file timestamps
showed.
\end{abstract}

\section*{Results at a glance}
\begin{center}
\renewcommand{\arraystretch}{1.15}
\begin{tabular}{@{}>{\raggedright}p{0.40\textwidth}>{\raggedright}p{0.34\textwidth}>{\raggedright\arraybackslash}p{0.20\textwidth}@{}}
\toprule
Situation & Bound & Where proved\\
\midrule
Any input & always sorted & Theorem~\ref{thm:correct}\\
Any input & $\le 5n\log_2 n+98n+C$ steps & Theorem~\ref{thm:worst}\\
Model exactly right (idealised) & $\le 34n+C$ expected & Corollary~\ref{cor:ideal}\\
Independent continuous data, model at $\chi^2$-distance $\chi^2$ & $\le(34+5\chi^2)n+O(\sqrt n)$ expected & Theorem~\ref{thm:main}\\
Uniform data, uniform model chosen & $\le 34n+O(\sqrt n)$ expected & Corollary~\ref{cor:uniform}\\
Any density with $\mu\le f\le M$ on a bounded interval, empirical model & $\le(29+15M/\mu)\,n+O(\sqrt n)$ expected & Corollary~\ref{cor:empirical}\\
A value with probability $\pi>0$ & $\ge\tfrac{\pi n}{2}\log_2\tfrac{\pi n}{2}$ expected, for any model & Proposition~\ref{prop:atoms}\\
\bottomrule
\end{tabular}
\end{center}
Here $C$ is a constant that does not depend on $n$ (it covers the work done
on the fixed-size sample). The constants are upper bounds chosen to make
each step easy to check, not predictions of the measured running time.

\paragraph{The idea in one paragraph.} Placing an element costs $O(1)$, so
placement is linear. Sorting a bucket of $k$ elements costs at most a
constant times $k^2$, so all the bucket sorts together cost
$O\bigl(\sum_j N_j^2\bigr)$, where $N_j$ is the size of bucket $j$. If the
model were the true CDF $F$, then $F(X)$ would be uniform on $[0,1]$
(Lemma~\ref{lem:pit}), so each element would land in each bucket with
probability exactly $1/n$. Each $N_j$ would then be
$\mathrm{Binomial}(n,1/n)$, with $\E[N_j^2]<2$, and the bucket sorts would
cost less than a constant times $2n$ in expectation. The rest of the proof
removes two idealisations: the model is only approximately $F$ (this costs
a factor $1+\chi^2$), and it is built from the same data it then sorts
(this costs only $O(\sqrt n)$).

\paragraph{Relation to the earlier proof.} \code{distribution\_sort\_proof.tex}
analyses an earlier version of the algorithm, which has recursion,
discrete models and an all-equal check. This document follows the same
structure for the current, simpler version. One thing is new here: the
earlier proof treated the model as independent of the data (its
assumption (A2)), whereas Lemma~\ref{lem:condition} below proves the
expected-time bound without that assumption.

%======================================================================
\section{The algorithm}\label{sec:alg}
%======================================================================

The code has two constants: \code{BIG\_BUCKET} $=64$ and
\code{SAMPLE\_SIZE} $=1000$. A call on a list of $n$ numbers performs these
steps.

\begin{enumerate}[label=\textbf{D\arabic*.},leftmargin=3em]
\item \textbf{Small input.} If $n\le 64$, insertion-sort a copy of the list
  and return it.
\item \textbf{Sample.} Let $s=\min(n,1000)$. Choose $s$ distinct positions
  uniformly at random with \code{random.sample}, and add the list's minimum
  $\xmin$ and maximum $\xmax$. Sort the resulting $s+2$ values into
  $q_0\le q_1\le\dots\le q_L$, where $L=s+1$. Note $q_0=\xmin$ and
  $q_L=\xmax$.
\item \textbf{Model.} Fit three candidates to $q_0,\dots,q_L$: a uniform CDF
  on $[q_0,q_L]$, a Normal CDF with the sample mean and standard deviation,
  and an Exponential CDF shifted to start at $q_0$. A candidate is accepted
  if its Kolmogorov--Smirnov distance to the sample is below
  $1.36/\sqrt{L+1}$; the accepted candidate with the smallest distance
  becomes the model $\Fh$. If none is accepted, $\Fh$ is the
  \emph{empirical CDF}: with $c(x)=\#\{i:q_i\le x\}$,
  \[
  \Fh(x)=
  \begin{cases}
  0 & c(x)=0,\\[2pt]
  \dfrac{1}{L}\Bigl(c(x)-1+\dfrac{x-q_{c(x)-1}}{q_{c(x)}-q_{c(x)-1}}\Bigr) & 1\le c(x)\le L,\\[8pt]
  1 & c(x)=L+1.
  \end{cases}
  \]
  (When $1\le c(x)\le L$ we have $q_{c(x)-1}\le x<q_{c(x)}$, so the
  denominator is positive.)
\item \textbf{Place.} Create $n$ empty buckets and append each element $x$
  to bucket
  \[
  \beta(x)=\min\bigl(n-1,\ \max(0,\ \mathrm{int}(n\,\Fh(x)))\bigr).
  \]
\item \textbf{Finish each bucket.} A bucket of size $k\le 64$ is
  insertion-sorted. A bucket of size $k>64$ is merge-sorted (the
  \emph{safety net}).
\item \textbf{Join} the buckets in order and return the result.
\end{enumerate}

\paragraph{Conventions.} We use the RAM model. A \emph{step} is one
iteration of any loop, including loops inside library routines (copying or
slicing a list, \code{min}, \code{max}, \code{sorted}, \code{bisect},
\code{random.sample}). The straight-line code around an iteration counts as
part of it; a call to \code{erf} or \code{exp} counts as one step; appending
to a list counts as one step (CPython over-allocates lists, so this holds in
amortised time up to a constant factor). Each step costs $O(1)$ machine
time. Real numbers are treated as exact (Remark~\ref{rem:float} discusses
floating point). All logarithms are base~2. $T(n)$ is the number of steps
of a call on $n$ elements.

%======================================================================
\section{Elementary facts}
%======================================================================

\begin{fact}[Insertion sort]\label{fact:ins}
Insertion sort on $k\ge1$ elements takes at most $(k-1)+\binom{k}{2}\le k^2$
steps.
\end{fact}
\begin{proof}
The outer loop runs $k-1$ times. Each inner-loop step moves one element
past the element being inserted, and each of the $\binom k2$ pairs is
moved past each other at most once. Finally
$(k-1)+\tfrac{k(k-1)}{2}\le k^2$.
\end{proof}

\begin{fact}[Merge sort]\label{fact:merge}
\code{merge\_sort} on $k\ge1$ elements takes at most $5k(\ceil{\log k}+1)$
steps and at least $k(\floor{\log k}+1)$ steps.
\end{fact}
\begin{proof}
A call on $k\ge2$ elements copies its input (\code{list(values)}), slices it
into halves of sizes $\floor{k/2}$ and $\ceil{k/2}$, recurses on both, and
merges. Using $\floor{\floor{x}/2}=\floor{x/2}$ and
$\ceil{\ceil{x}/2}=\ceil{x/2}$, every call at recursion depth $d$ has
between $\floor{k/2^d}$ and $\ceil{k/2^d}$ elements.

\emph{Upper bound.} At depth $\ceil{\log k}$ every call has one element, so
there are at most $\ceil{\log k}+1$ depths. The calls at one depth work on
disjoint parts of the input. Per element at one depth, a call spends at most
5 steps: one to copy it, one to slice it, at most two in merging (one loop
iteration, or one slice plus one extend for leftovers), and one for its
share of call overhead (a depth has at most as many calls as elements).

\emph{Lower bound.} For $d\le\floor{\log k}$ every call at depth $d$ has at
least $\floor{k/2^d}\ge1$ element, and every call at depth $d-1$ has at
least $2$ elements and therefore recurses. So every element is in some call
at each depth $0,1,\dots,\floor{\log k}$, and each such call copies it once.
\end{proof}

\begin{fact}[Binary search]\label{fact:bs}
Evaluating the empirical CDF uses \code{bisect\_right} on $L+1\le 1002$
points, which takes at most $10$ iterations.
\end{fact}
\begin{proof}
The search interval starts with length $L+1\le1002$ and each iteration at
least halves it (rounding down), so there are at most
$\floor{\log 1002}+1=10$ iterations.
\end{proof}

\begin{fact}[\code{random.sample}]\label{fact:sample}
Drawing $s\le n$ positions with \code{random.sample} takes at most $5n$
steps in expectation, and the number of steps is independent of which
positions are returned.
\end{fact}
\begin{proof}
CPython computes $\sigma=21$, plus $4^{\ceil{\log_4(3s)}}\ge3s$ if $s>5$.
Each random draw calls \code{randbelow}, a rejection loop that takes fewer
than 2 iterations on average.
\begin{itemize}
\item If $n\le\sigma$, it allocates the result ($s$ steps), copies the list
  ($n$ steps) and makes $s$ draws and swaps (fewer than $3s$ steps on
  average), at most $n+4s\le5n$ in total.
\item Otherwise $n>\sigma\ge3s$, and it draws positions until it has $s$
  distinct ones. While fewer than $s$ have been chosen, a draw repeats an
  earlier one with probability below $s/n<1/3$, so each new position needs
  fewer than $1.5$ draws on average. The total is at most
  $s+1.5s\cdot3+s<7s<3n$ on average.
\end{itemize}
In both cases the only randomness in the step count comes from rejection
loops, and in a rejection loop the number of trials is independent of the
value finally accepted. So the step count is independent of the positions
returned.
\end{proof}

\begin{fact}[Second moment of a binomial]\label{fact:binom}
If $N\sim\mathrm{Binomial}(n,p)$, then $\E[N^2]=np(1-p)+n^2p^2\le np+n^2p^2$.
\end{fact}

\begin{fact}[Cauchy--Schwarz inequality]\label{fact:cs}
For non-negative functions $u,v$ on an interval $I$,
$\int_I uv\le\bigl(\int_I u^2\bigr)^{1/2}\bigl(\int_I v^2\bigr)^{1/2}$.
\end{fact}

%======================================================================
\section{Correctness}
%======================================================================

\begin{lemma}[The model never decreases]\label{lem:mono}
Every model built in step D3 is non-decreasing on $[\xmin,\xmax]$, which
contains every input value. Hence the bucket map is non-decreasing: if
$x\le y$ then $\beta(x)\le\beta(y)$.
\end{lemma}
\begin{proof}
The Normal model is a CDF. The uniform model $(x-q_0)/(q_L-q_0)$ is
increasing. The Exponential model is $0$ up to $q_0$ and
$1-e^{-(x-q_0)/\lambda}$ after it, with $\lambda>0$; both pieces are
non-decreasing and they agree at $q_0$.

For the empirical CDF, $c(x)$ is non-decreasing in $x$. While $c(x)$ stays
equal to some $c\in[1,L]$, the formula is linear with positive slope and
takes values in $[\frac{c-1}{L},\frac{c}{L})$. When $c(x)$ jumps from $c$ to
$c'>c$ (at a sample point, possibly repeated), the value just before the
jump is below $\frac cL$ and the value at the jump is $\frac{c'-1}{L}\ge\frac cL$,
or $1$ if $c'=L+1$. Below $q_0$ the value is $0$, the smallest possible.

Finally, $\beta$ composes $\Fh$ with three non-decreasing maps: multiplying
by $n$, \code{int}, and clamping to $[0,n-1]$. (\code{int} truncates towards
zero rather than rounding down, but the two differ only on negative
numbers, which the clamp sends to $0$ either way.)
\end{proof}

\begin{theorem}[Correctness]\label{thm:correct}
For every input list, \code{distribution\_sort} stops and returns its
elements in non-decreasing order.
\end{theorem}
\begin{proof}
If $n\le64$, insertion sort is correct. Otherwise, let $x$ be in bucket
$j$ and $y$ in bucket $j'$ with $j<j'$. If $x>y$, Lemma~\ref{lem:mono} would
give $j=\beta(x)\ge\beta(y)=j'$, a contradiction. So every element of an
earlier bucket is at most every element of a later bucket. Each bucket is
sorted by insertion sort or merge sort, both of which are correct and stop.
Joining sorted buckets that are in order gives a sorted list.
\end{proof}

%======================================================================
\section{Where the time goes}
%======================================================================

\begin{lemma}[Work outside the bucket sorts]\label{lem:overhead}
For $n>64$, the work outside the bucket sorts in step D5 is $W\le A\,n+C$
in expectation, with $A=24$ and $C$ a constant that does not depend on
$n$. The bound also holds conditional on the positions sampled and on the
data.
\end{lemma}
\begin{proof}
The table lists every part of steps D1--D6 except the bucket sorts.
\begin{center}
\renewcommand{\arraystretch}{1.15}
\begin{tabular}{@{}>{\raggedright}p{0.36\textwidth}r>{\raggedright\arraybackslash}p{0.46\textwidth}@{}}
\toprule
Work & Steps $\le$ & Reason\\
\midrule
D1 copy the input & $n$ & one pass\\
D2 \code{random.sample} & $5n$ & Fact~\ref{fact:sample}, in expectation\\
D2 $\min$ and $\max$ & $2n$ & one pass each\\
D2 build the sample list & $n+2$ & $s+2\le n+2$ elements\\
D2 sort the sample; D3 fit three models, KS distances, build $\Fh$ & $C$ & every loop runs over the $L+1\le1002$ sample points, so the work is bounded by a constant\\
D4 create $n$ buckets & $n$ & one pass\\
D4 place each element & $12n$ & per element: 1 iteration, $\le10$ search steps (Fact~\ref{fact:bs}), 1 append\\
D5, D6 visit buckets and join & $2n$ & one iteration per bucket, one copy per element\\
\midrule
Total & $24n+C$ &\\
\bottomrule
\end{tabular}
\end{center}
The constant $2$ from the sample list is absorbed into $C$. Only the
\code{random.sample} row is random, and by Fact~\ref{fact:sample} its step
count is independent of the positions returned and of the data (which the
sampler never looks at), so conditioning on them does not change its
expectation.
\end{proof}

\begin{lemma}[Cost of one bucket]\label{lem:bucket}
Let $\mathrm{cost}(k)$ be the steps used to sort a bucket of size $k$ in
step D5. Then
\[
\mathrm{cost}(k)\le
\begin{cases}
k^2\le 64k & k\le 64,\\
5k(\ceil{\log k}+1) & k>64,
\end{cases}
\qquad\text{and in all cases}\qquad \mathrm{cost}(k)\le 5k^2 .
\]
\end{lemma}
\begin{proof}
The first case is Fact~\ref{fact:ins} and the second is
Fact~\ref{fact:merge}. Since $\ceil{\log k}+1\le k$ for every $k\ge1$,
both cases are at most $5k^2$.
\end{proof}

Writing $N_j$ for the size of bucket $j$, the two lemmas give, for $n>64$,
\begin{equation}\label{eq:split}
T(n)=W+\sum_{j=0}^{n-1}\mathrm{cost}(N_j),\qquad \E[W]\le An+C,\qquad
\sum_j N_j=n .
\end{equation}

%======================================================================
\section{Worst case: \texorpdfstring{$O(n\log n)$}{O(n log n)} for every input}
%======================================================================

\begin{theorem}[Worst case]\label{thm:worst}
For every input with $n>64$ elements,
\[
\E[T(n)]\le 5n\log n+(A+74)\,n+C=5n\log n+98n+C,
\]
where the expectation is only over \code{random.sample}'s internal draws.
For $n\le64$, $T(n)\le n+n^2\le 65n$.
\end{theorem}
\begin{proof}
By Lemma~\ref{lem:bucket}, a bucket with $N_j\le64$ costs at most $64N_j$,
and a bucket with $N_j>64$ costs at most
$5N_j(\ceil{\log N_j}+1)\le5N_j(\log n+2)$, since $N_j\le n$. Adding over
all buckets and using $\sum_j N_j=n$,
\[
\sum_j\mathrm{cost}(N_j)\le 64n+5n(\log n+2)=5n\log n+74n .
\]
Add $\E[W]\le An+C$ from \eqref{eq:split}. For $n\le64$, step D1 copies the
list and insertion-sorts it (Fact~\ref{fact:ins}).
\end{proof}

\begin{remark}
This bound is what the safety net buys. Without it, one overfull bucket of
size $\Theta(n)$ would be insertion-sorted in $\Theta(n^2)$ time. Replacing
\code{random.sample} by a partial Fisher--Yates shuffle would make the
bound hold on every run, without the expectation.
\end{remark}

%======================================================================
\section{Expected linear time for continuous data}
%======================================================================

\subsection{The ideal case}

\begin{lemma}[Probability integral transform]\label{lem:pit}
If $X$ has a continuous CDF $F$, then $F(X)$ is uniform on $[0,1]$.
\end{lemma}
\begin{proof}
Fix $u\in(0,1)$. Since $F$ is continuous and non-decreasing, with limits $0$
and $1$, the set $\{x:F(x)\le u\}$ is non-empty and bounded above, and it
contains its supremum $x_u$; continuity also gives $F(x_u)=u$. Then
$F(X)\le u$ exactly when $X\le x_u$, so
$\Prob(F(X)\le u)=F(x_u)=u$.
\end{proof}

\begin{corollary}[Ideal model]\label{cor:ideal}
Suppose the data are independent draws from a continuous CDF $F$ and,
hypothetically, $\Fh=F$ and $\Fh$ were chosen independently of the data.
Then $\E[T(n)]\le(A+10)\,n+C=34n+C$.
\end{corollary}
\begin{proof}
By Lemma~\ref{lem:pit}, element $i$ lands in bucket $j$ when $nF(X_i)$ lies in
$[j,j+1)$ (or in $[n-1,n]$ for the last bucket), which has probability
exactly $1/n$. The elements are independent, so
$N_j\sim\mathrm{Binomial}(n,1/n)$ and, by Fact~\ref{fact:binom},
$\E[N_j^2]=1-\tfrac1n+1<2$. By Lemma~\ref{lem:bucket} and
\eqref{eq:split},
$\E[T]\le An+C+5\sum_j\E[N_j^2]<An+C+10n$.
\end{proof}

The code differs from this ideal in two ways: $\Fh$ is only an
approximation of $F$, and $\Fh$ is computed from a sample of the same list
it then sorts. The next two subsections deal with each.

\subsection{The model comes from the data}

\paragraph{Assumption (A).} The input values $X_1,\dots,X_n$ are independent
and identically distributed with a density $f$ and CDF $F$. (So, with
probability 1, no two values are equal.)

\paragraph{Notation.} Let $J$ be the set of sampled positions, $\iota$ and
$\kappa$ the positions of the minimum $\xmin$ and maximum $\xmax$, and
$E=J\cup\{\iota,\kappa\}$, so $|E|\le s+2$. Let $R$ be the remaining
positions. Let $\G$ denote the information $(J,\iota,\kappa,(X_i)_{i\in E})$.
Write $Z=F(\xmax)-F(\xmin)$ and let
\[
f_\ast(x)=\frac{f(x)}{Z}\quad\text{for }\xmin<x<\xmax,\qquad f_\ast(x)=0\text{ otherwise}
\]
be the density $f$ restricted to $(\xmin,\xmax)$.

\begin{lemma}[Conditioning on the model's inputs]\label{lem:condition}
\textup{(i)} The model $\Fh$, and the bucket of every element in $E$, are
determined by $\G$.
\textup{(ii)} Given $\G$, the values $(X_i)_{i\in R}$ are independent with
density $f_\ast$.
\end{lemma}
\begin{proof}
(i) Steps D2 and D3 use only the sampled values, $\xmin=X_\iota$ and
$\xmax=X_\kappa$.

(ii) \code{random.sample} uses its own random numbers, so $J$ is independent
of the data. Fix $J$, $\iota=i_0$ and $\kappa=k_0$ (with $i_0\ne k_0$, which
holds with probability 1 when $n\ge2$). The joint density of the data is
$\prod_i f(x_i)$, restricted to the event that $x_{i_0}<x_i<x_{k_0}$ for
every $i\ne i_0,k_0$. Now also fix the values $x_i$ for $i\in E$. The
conditions involving positions in $E$ are then fixed, and what remains of
the density of $(x_i)_{i\in R}$ is proportional to
$\prod_{i\in R}f(x_i)\,\mathbf 1\{\xmin<x_i<\xmax\}$. This is a product, so
the values are independent, each with density proportional to
$f\cdot\mathbf 1_{(\xmin,\xmax)}$, which is $f_\ast$.
\end{proof}

\subsection{The model is only approximately right}

From here on, condition on $\G$, so that $\Fh$ is a fixed function. Every
model in D3 is continuous with a derivative $\fh$ except at finitely many
points, apart from upward jumps of the empirical CDF at repeated sample
points (which happen with probability 0 under Assumption~(A)).

\begin{definition}
The \emph{$\chi^2$-distance} from the data to the model is
\[
\chi^2=\chi^2(f_\ast\,\|\,\fh)=\int_{\xmin}^{\xmax}\frac{f_\ast(x)^2}{\fh(x)}\,dx-1,
\]
taken to be $+\infty$ if $\fh=0$ on a set where $f_\ast>0$. By
Fact~\ref{fact:cs}, $1=\int f_\ast\le(\int f_\ast^2/\fh)^{1/2}(\int\fh)^{1/2}
\le(\int f_\ast^2/\fh)^{1/2}$, so $\chi^2\ge0$, with equality exactly when
$\fh=f_\ast$. For each bucket $j$ let
$I_j=\{x\in(\xmin,\xmax):\beta(x)=j\}$ and $p_j=\int_{I_j}f_\ast$.
\end{definition}

\begin{lemma}[The model gives each bucket mass at most $1/n$]\label{lem:mass}
For every bucket $j$, $\int_{I_j}\fh\le1/n$.
\end{lemma}
\begin{proof}
By Lemma~\ref{lem:mono}, $I_j$ is an interval. On $I_j$, $n\Fh$ lies in
$[j,j+1)$, or in $[0,1)$ for $j=0$, or in $[n-1,n]$ for $j=n-1$ (because of
the clamp, using $0\le\Fh\le1$ on $[\xmin,\xmax]$). So $\Fh$ increases by at
most $1/n$ across $I_j$. Split $I_j$ at the finitely many points where
$\Fh$ is not differentiable. On each piece, $\int\fh$ equals the increase of
$\Fh$ across the piece, and jumps at the split points only add to the
total increase. So $\int_{I_j}\fh\le1/n$.
\end{proof}

\begin{lemma}[Collisions are rare]\label{lem:collide}
$\displaystyle\sum_j p_j^2\le\frac{1+\chi^2}{n}$.
\end{lemma}
\begin{proof}
If $\chi^2=\infty$ there is nothing to prove. Otherwise, by Fact~\ref{fact:cs}
and Lemma~\ref{lem:mass},
\[
p_j=\int_{I_j}\frac{f_\ast}{\sqrt{\fh}}\cdot\sqrt{\fh}
\le\Bigl(\int_{I_j}\frac{f_\ast^2}{\fh}\Bigr)^{1/2}\Bigl(\int_{I_j}\fh\Bigr)^{1/2}
\le\Bigl(\frac1n\int_{I_j}\frac{f_\ast^2}{\fh}\Bigr)^{1/2}.
\]
Square, and add over $j$; the $I_j$ do not overlap and cover
$(\xmin,\xmax)$, so $\sum_jp_j^2\le\frac1n\int f_\ast^2/\fh=\frac{1+\chi^2}{n}$.
\end{proof}

\begin{lemma}[Bucket sizes]\label{lem:moments}
$\displaystyle\E\Bigl[\sum_j N_j^2\,\Big|\,\G\Bigr]\le(2+\chi^2)\,n+2(s+2)\sqrt{(1+\chi^2)\,n}+(s+2)^2.$
\end{lemma}
\begin{proof}
Split $N_j=N_j'+S_j$, where $N_j'$ counts elements of $R$ in bucket $j$ and
$S_j$ counts elements of $E$. By Lemma~\ref{lem:condition}, the $S_j$ are
fixed given $\G$, with $\sum_jS_j=|E|\le s+2$, and each element of $R$
lands in bucket $j$ independently with probability $p_j$, so
$N_j'\sim\mathrm{Binomial}(|R|,p_j)$ with $|R|\le n$. Expanding
$N_j^2=N_j'^2+2S_jN_j'+S_j^2$:
\begin{itemize}
\item By Fact~\ref{fact:binom}, $\sum_j\E[N_j'^2]\le\sum_j(np_j+n^2p_j^2)
  \le n+n^2\cdot\frac{1+\chi^2}{n}=(2+\chi^2)\,n$, using $\sum_jp_j=1$ and
  Lemma~\ref{lem:collide}.
\item $\sum_j2S_j\E[N_j']\le2n\,(\max_jp_j)\sum_jS_j\le2n(s+2)\bigl(\sum_jp_j^2\bigr)^{1/2}
  \le2(s+2)\sqrt{(1+\chi^2)\,n}$.
\item $\sum_jS_j^2\le\bigl(\sum_jS_j\bigr)^2\le(s+2)^2$.\qedhere
\end{itemize}
\end{proof}

\subsection{The main theorem}

\begin{theorem}[Expected linear time]\label{thm:main}
Under Assumption~(A), for $n>64$,
\[
\E[T(n)\mid\G]\le(A+10+5\chi^2)\,n+10(s+2)\sqrt{(1+\chi^2)\,n}+5(s+2)^2+C .
\]
With $A=24$ and $s+2\le1002$, this is $(34+5\chi^2)\,n+O(\sqrt n)$. Averaging
over $\G$ (and using $\E\sqrt Y\le\sqrt{\E Y}$),
\[
\E[T(n)]\le\bigl(34+5\,\E[\chi^2]\bigr)\,n+10(s+2)\sqrt{(1+\E[\chi^2])\,n}+5(s+2)^2+C ,
\]
so the expected time is linear whenever $\E[\chi^2]$ is bounded
independently of $n$. Theorem~\ref{thm:worst} applies in every case.
\end{theorem}
\begin{proof}
By \eqref{eq:split} and Lemma~\ref{lem:bucket},
$T\le W+5\sum_jN_j^2$. By Lemma~\ref{lem:overhead}, $\E[W\mid\G]\le An+C$.
Multiply the bound of Lemma~\ref{lem:moments} by 5 and add.
\end{proof}

\begin{remark}[Why the total is linear]
The proof charges every crowded bucket a \emph{quadratic} cost, yet the
total is still linear. The reason is Lemma~\ref{lem:collide}: two random
elements share a bucket with probability only about $(1+\chi^2)/n$, so
crowded buckets are rare enough that even a quadratic cost adds up to
$O(n)$. The constant $5(s+2)^2\approx5\times10^6$ comes from the at most
$1002$ sampled points possibly sharing buckets. It is a crude bound, but it
does not grow with $n$.
\end{remark}

\subsection{When is \texorpdfstring{$\chi^2$}{chi-squared} small?}

\begin{corollary}[Uniform data]\label{cor:uniform}
If the data are uniform on an interval and the uniform model is chosen,
then $\chi^2=0$ and
$\E[T(n)\mid\G]\le34n+10(s+2)\sqrt n+5(s+2)^2+C$.
\end{corollary}
\begin{proof}
The uniform model is $\Fh(x)=(x-\xmin)/(\xmax-\xmin)$, so
$\fh=1/(\xmax-\xmin)$ on $(\xmin,\xmax)$. A uniform density restricted to
$(\xmin,\xmax)$ and rescaled is again uniform there, so $f_\ast=\fh$ and
$\chi^2=0$. Apply Theorem~\ref{thm:main}.
\end{proof}

\begin{corollary}[Empirical model]\label{cor:empirical}
Suppose $f$ is supported on a bounded interval $[\alpha,\beta]$, with
$0<\mu\le f\le M$ there, and consider the empirical model of step D3. Then
\[
\E\bigl[1+\chi^2\,\big|\,J,\iota,\kappa,\xmin,\xmax\bigr]\le\frac{2(s+1)}{s}\cdot\frac M\mu\le\frac{3M}{\mu}.
\]
Hence the variant of the algorithm that always uses the empirical model
satisfies $\E[T(n)]\le(29+15M/\mu)\,n+O(\sqrt n)$.
\end{corollary}
\begin{proof}
Let $s'=|J\setminus\{\iota,\kappa\}|\ge s-2$ be the number of sampled values
strictly between $\xmin$ and $\xmax$. By the argument of
Lemma~\ref{lem:condition}(ii), given $(J,\iota,\kappa,\xmin,\xmax)$ these
$s'$ values are independent with density $f_\ast$. Mapped through the CDF
$F_\ast$ of $f_\ast$, they become $s'$ independent uniform points on
$[0,1]$, while $\xmin\mapsto0$ and $\xmax\mapsto1$.

The sample points split $(\xmin,\xmax)$ into $s'+1$ segments of positive
length. (A repeated point, which occurs when $J$ contains $\iota$ or
$\kappa$, gives a segment of length zero, which does not matter.) On a
segment of length $\ell_k$, the empirical CDF rises by exactly $1/L$, so
$\fh=1/(L\ell_k)$. Let $\omega_k=\int_{\text{segment }k}f_\ast$. On
$(\xmin,\xmax)$ we have $\mu/Z\le f_\ast\le M/Z$, so
$\omega_k\ge\mu\ell_k/Z$, that is $\ell_k\le Z\omega_k/\mu$, and
$\int_{\text{segment }k}f_\ast^2\le(M/Z)\,\omega_k$. Therefore
\[
1+\chi^2=\sum_kL\ell_k\int_{\text{segment }k}f_\ast^2\le\sum_kL\cdot\frac{Z\omega_k}{\mu}\cdot\frac{M}{Z}\,\omega_k
=\frac{LM}{\mu}\sum_k\omega_k^2 .
\]
The $\omega_k$ are the gaps between $s'$ independent uniform points on
$[0,1]$ (including the gaps to $0$ and to $1$). Each of the $s'+1$ gaps has
the $\mathrm{Beta}(1,s')$ distribution, with second moment
$\frac{2}{(s'+1)(s'+2)}$, so $\E\sum_k\omega_k^2=\frac{2}{s'+2}$. With
$L=s+1$ and $s'+2\ge s$, this gives the first bound. The second uses
$s=\min(n,1000)\ge65$.

For the running time, apply Theorem~\ref{thm:main} with
$\E[\chi^2]\le3M/\mu-1$: $34+5(3M/\mu-1)=29+15M/\mu$.
\end{proof}

In the actual code, the empirical model is replaced by a named model only
when the named model passes the Kolmogorov--Smirnov test, and in that case
Theorem~\ref{thm:main} applies with the named model's $\chi^2$.

\begin{remark}[Normal data with a fitted Normal model]
For a Normal density $f$ with standard deviation $\sigma$ and a Normal model
with $\hat\sigma$, the integral $\int f^2/\fh$ is a Gaussian integral whose
exponent has $x^2$-coefficient $-\frac1{\sigma^2}+\frac1{2\hat\sigma^2}$. So
it is finite exactly when $\hat\sigma^2>\sigma^2/2$. For small fitting
errors, a second-order expansion (with the Fisher information as
coefficients) gives
$\chi^2(f\,\|\,\fh)\approx(\Delta\mu/\sigma)^2+2(\Delta\sigma/\sigma)^2$.
With about 1000 sample points the relative errors are about $3\%$ and
$2\%$, so $\chi^2\approx0.002$. Restricting to $(\xmin,\xmax)$ changes little:
$1+\chi^2(f_\ast\|\fh)\le(1+\chi^2(f\|\fh))/Z^2$, and $Z\to1$ as $n$ grows.
\end{remark}

%======================================================================
\section{Where linearity fails: repeated values}
%======================================================================

\begin{proposition}[A repeated value forces $\Omega(n\log n)$]\label{prop:atoms}
Suppose the values are independent and some value $v$ has
$\Prob(X=v)=\pi>0$. Then for every $n\ge130/\pi$, whatever model step D3
chooses,
\[
\E[T(n)]\ge\bigl(1-e^{-\pi n/8}\bigr)\,\frac{\pi n}{2}\,\log\frac{\pi n}{2}.
\]
\end{proposition}
\begin{proof}
Let $K$ be the number of copies of $v$, so $K\sim\mathrm{Binomial}(n,\pi)$.
All copies have the same bucket $\beta(v)$, because $\beta$ is a function
of the value: no model can separate equal values. By the Chernoff bound,
$\Prob(K<\pi n/2)\le e^{-\pi n/8}$. If $K\ge\pi n/2\ge65$, that bucket has
more than 64 elements and is merge-sorted, which by Fact~\ref{fact:merge}
costs at least $N(\floor{\log N}+1)\ge K\log K\ge\frac{\pi n}2\log\frac{\pi n}2$,
where $N\ge K$ is the bucket's size and $x\log x$ is increasing.
\end{proof}

\begin{remark}[The benchmarks, and the fix]
This is exactly the timestamp benchmark: $73\%$ of the files share one
modification time, so $\pi\approx0.73$, and the sort was slower than plain
merge sort. File sizes have many smaller repeated values. (A density does
not exist for such data, so $\chi^2$ is undefined.)

The proof also shows what to change. The expensive step is merge-sorting a
bucket whose values are all equal. A single pass that checks whether all
values in a bucket are equal costs $O(N)$ and makes such a bucket free to
finish, which removes the $\Omega(n\log n)$ term for buckets holding one
repeated value. Buckets that mix a repeated value with other values need a
separate argument, like the pure/mixed bucket analysis in the earlier
proof.
\end{remark}

%======================================================================
\section{Discussion}
%======================================================================

\begin{remark}[Why the $\Omega(n\log n)$ sorting lower bound does not apply]
Any \emph{comparison} sort needs $\Omega(n\log n)$ comparisons, even on
average over random inputs, because it must tell apart $n!$ orderings. Step
D4 does not compare elements with each other: it computes $\floor{n\Fh(x)}$
with arithmetic, which locates $x$ among $n$ buckets in $O(1)$ time. All
of the speed comes from this, and it pays off only when the model fits
($\chi^2$ bounded).
\end{remark}

\begin{remark}[Why the sample size must be a constant]
The empirical CDF finds $x$'s segment by binary search over the sample,
which costs $\log(\text{sample size})$ steps. With a fixed sample of 1000
this is at most 10 steps (Fact~\ref{fact:bs}). If the sample grew with $n$,
placement alone would cost $\Theta(n\log n)$. The price of a fixed sample
is that the empirical model's accuracy, and so its $\chi^2$, does not
improve as $n$ grows; Corollary~\ref{cor:empirical} shows that it stays
bounded anyway.
\end{remark}

\begin{remark}[Floating point]\label{rem:float}
The proofs assume exact arithmetic. In floating point, correctness depends
only on Lemma~\ref{lem:mono}, that is, on the computed $\Fh$ being
non-decreasing. This holds because IEEE subtraction and division by a
positive number are monotone, \code{erf} and \code{exp} are monotone in
standard maths libraries, and the empirical CDF's interpolation fraction
stays in $[0,1]$ after rounding. It also needs differences such as
$\xmax-\xmin$ not to overflow, which holds for values below about
$10^{308}$ in size. The running-time bounds do not depend on rounding.
\end{remark}

\begin{remark}[Connection to the measurements]
For uniform and Normal data the fitted models have $\chi^2\approx0$, and the
measured time per element stayed roughly flat from $10^4$ to $10^6$
elements, as Theorem~\ref{thm:main} predicts. The extra slowdown observed
above $10^5$ elements also hit merge sort and Python's built-in sort, so it
comes from the machine rather than the algorithm (part of it was measured
to be the garbage collector; the rest is most likely memory-cache misses),
and the RAM model does not capture it.
\end{remark}

\appendix
\section{Where every number comes from}
\renewcommand{\arraystretch}{1.2}
\begin{longtable}{@{}>{\raggedright}p{0.17\textwidth}p{0.77\textwidth}@{}}
\toprule
Number & Origin\\
\midrule
\endhead
$64$ & \code{BIG\_BUCKET} in the code: the largest bucket that is insertion-sorted.\\
$1000$ & \code{SAMPLE\_SIZE} in the code. With $\xmin$ and $\xmax$ added, the model sees at most $1002$ points.\\
$1.36$ & KS acceptance threshold factor: the $95\%$ point of the Kolmogorov distribution. It affects which model is chosen, not the bounds.\\
$10$ & Binary-search iterations on $\le1002$ points: $\floor{\log1002}+1=10$ (Fact~\ref{fact:bs}).\\
$5$ & Steps per element per merge-sort level: copy, slice, two for merging, call overhead (Fact~\ref{fact:merge}). Also $\mathrm{cost}(k)\le5k^2$ (Lemma~\ref{lem:bucket}).\\
$5n$ & Expected steps of \code{random.sample} (Fact~\ref{fact:sample}).\\
$24$ ($A$) & Sum of the per-element costs in Lemma~\ref{lem:overhead}.\\
$74$ & $64$ (small buckets) $+\,5\cdot2$ (from $\ceil{\log k}+1\le\log n+2$) in Theorem~\ref{thm:worst}.\\
$98$ & $A+74$.\\
$10$, $34$ & $5\cdot2$ from $\E[N_j^2]<2$; and $A+10$ (Corollary~\ref{cor:ideal}, Theorem~\ref{thm:main}).\\
$2+\chi^2$ & $\sum_j\E[N_j'^2]\le n+n^2\cdot\frac{1+\chi^2}{n}$ (Lemma~\ref{lem:moments}).\\
$s+2$ & At most this many elements (the sample, $\xmin$, $\xmax$) are not independent of the model (Lemma~\ref{lem:condition}).\\
$2/(s'+2)$ & Expected sum of squared gaps between $s'$ uniform points (Corollary~\ref{cor:empirical}).\\
$3M/\mu$ & $2(s+1)/s\cdot M/\mu$ with $s\ge65$ (Corollary~\ref{cor:empirical}).\\
$29$, $15$ & $34+5(3M/\mu-1)$ (Corollary~\ref{cor:empirical}).\\
$130/\pi$, $8$ & Need $\pi n/2\ge65$ for the safety net to trigger; Chernoff bound $\Prob(K<\tfrac12\E K)\le e^{-\E K/8}$ (Proposition~\ref{prop:atoms}).\\
\bottomrule
\end{longtable}

\end{document}
